% ============================================
% 共享 Preamble - 适用于所有增强版 TikZ 图表
% 用于 NeurIPS, CVPR, ICLR, IEEE TPAMI 等顶会/顶刊
% ============================================

% 数学符号支持
\usepackage{amssymb, amsmath}

% 字体配置（支持中文）
\usepackage{fontspec}
\usepackage{xeCJK}
\setCJKmainfont{SimHei}  % 中文黑体（可根据系统调整）

% TikZ 核心包
\usepackage{tikz}
\usepackage{pgfplots}
\pgfplotsset{compat=1.18}

% TikZ 必备库
\usetikzlibrary{
  arrows.meta,          % 箭头样式
  positioning,          % 相对定位（禁止绝对坐标）
  shapes.geometric,     % 几何形状
  calc,                 % 坐标计算
  backgrounds,          % 背景层（用于 fit）
  fit,                  % 节点分组
  decorations.pathreplacing,  % 装饰（大括号等）
  patterns,             % 填充图案
  shapes.symbols        % 符号形状
}

% ============================================
% 配色方案：低饱和度马卡龙（禁止纯色）
% ============================================
\definecolor{mBlue}{HTML}{AEC6CF}      % 编码端/IR层/高层算子
\definecolor{mPink}{HTML}{FFD1DC}      % Discriminator/Baseline
\definecolor{mGreen}{HTML}{B1D8B7}     % 信道/调优层/底层算子
\definecolor{mPurple}{HTML}{B39EB5}    % 解码端/代码生成/Interpreter
\definecolor{mYellow}{HTML}{FDFD96}    % 半自动搜索/特殊标注
\definecolor{mGray}{HTML}{E5E4E2}      % 硬件/IO/中性背景
\definecolor{mDark}{HTML}{4A4A4A}      % 文字/线条（替代纯黑）

% ============================================
% 全局 TikZ 样式配置
% ============================================
\tikzset{
  % 默认箭头与线条
  >=stealth,
  thick,
  draw=mDark,

  % 基础节点样式
  block/.style={
    rectangle,
    rounded corners=3pt,
    draw=mDark,
    align=center,
    font=\sffamily\small,
    minimum height=0.85cm,
    inner sep=5pt
  },

  % 微观模块样式（拆解黑盒用）
  micro/.style={
    rectangle,
    rounded corners=2pt,
    draw=mDark!70,
    align=center,
    font=\sffamily\tiny,
    minimum height=0.4cm,
    inner sep=3pt
  },

  % 数学操作符节点（⊕加 ⊗乘 ‖拼接）
  op/.style={
    circle,
    draw=mDark,
    minimum size=0.45cm,
    inner sep=0,
    font=\scriptsize,
    fill=mGray!60
  },

  % 张量维度标注（信息密度提升）
  ts/.style={
    font=\tiny\sffamily,
    color=mDark!80
  },

  % 箭头样式
  arr/.style={
    -{Stealth[length=4pt]},
    thick,
    draw=mDark
  },

  % 虚线箭头（用于训练路径或可选路径）
  darr/.style={
    -{Stealth[length=4pt]},
    thick,
    draw=mDark,
    dashed
  },

  % 图例框样式
  legendbox/.style={
    rectangle,
    draw=mDark!50,
    rounded corners=3pt,
    fill=mGray!15,
    inner sep=5pt,
    font=\tiny\sffamily
  },
}

% ============================================
% pgfplots 默认样式（用于柱状图/折线图）
% ============================================
\pgfplotsset{
  every axis/.append style={
    font=\sffamily\small,
    tick label style={font=\sffamily\scriptsize, color=mDark},
    label style={font=\sffamily\small, color=mDark},
    legend style={
      font=\sffamily\tiny,
      draw=mDark!50,
      fill=mGray!15,
      rounded corners=2pt
    },
    grid style={draw=mGray, thin},
  }
}

% ============================================
% 注意事项
% ============================================
% 1. 禁止使用绝对坐标 `\node at (1,2)`，必须使用 `positioning` 库
% 2. 禁止使用 TikZ 默认纯色（red, blue），必须使用 mBlue, mGreen 等
% 3. 连线转折必须使用正交连线（`|-` 或 `-|`），禁止斜线乱飞
% 4. 每张图必须补充张量维度、微观结构、数学公式等信息密度元素
% 5. 使用 `fit` + `backgrounds` 实现层次化分组，体现模块归属关系
